\section{Conclusiones}
Las cinco resonancias registradas se asociaron con los armónicos impares $1,3,5,7,9$, en concordancia con las condiciones de frontera de un tubo cerrado por un extremo. El ajuste dio $R^2=0.9987$ y un intercepto compatible con cero en la escala de su error estándar.

La estimación por armónicos fue $v_h=(341.3\pm7.7)$ m/s, conservando la incerteza del registro original, frente a la referencia térmica $(345.9\pm0.6)$ m/s. La discrepancia de 1.32\% indica concordancia en la escala declarada, sin demostrar precisión o reproducibilidad por sí sola.

La reducción de longitud efectiva de 37.50 a 14.58 cm elevó la frecuencia de 218.23 a 596.73 Hz. El ajuste frente a longitud inversa produjo 361.3 m/s, con error estándar estadístico de 4.3 m/s y discrepancia de 4.45\%. Su intercepto de $-24.2$ Hz evidencia una desviación del modelo ideal; la mayor cercanía del ajuste al origen no justifica ocultar esa discrepancia.

Los registros permiten abordar la generación de resonancias, su asociación modal y las estimaciones de velocidad y dependencia con longitud. Los volúmenes teóricos para diapasones se calcularon, pero su comprobación experimental no pudo verificarse. Fijar la geometría del montaje, repetir barridos y registrar volúmenes finales son mejoras prioritarias para evaluar sesgos y reproducibilidad.
